% GENERATED FILE. Edit canonical lessons, then run npm run generate:books.
% canonical-source-hash: fnv1a64:0c0aa10281bfda89
% canonical-lessons: ES-C708-celos, ES-C708-celoso, ES-C708-envidia, ES-C708-envidiar, ES-C708-rencor

\chapter{Celos y envidia}
\label{ch:es-celos-y-envidia}
\hlchaptermodality{708}

\begin{chapteropening}
\textbf{By the end of this chapter:} \emph{I can talk about jealousy and envy between friends and partners.}

Complete the last lesson of chapter 708: I can talk about jealousy and envy between friends and partners.
\end{chapteropening}

\section[los celos]{los celos — jealousy}

\label{lesson:ES-C708-celos}



\begin{quote}
Before the new one: say the Spanish for anger, then the Spanish for to annoy.
\end{quote}

\subsection*{You'll want to know: los celos}
\textbf{los celos} — \textquotedblleft{}jealousy\textquotedblright{}.

\begin{quote}
Los celos destruyen una pareja.
\end{quote}

\begin{quote}
\emph{Jealousy destroys a couple.}
\end{quote}

Plural in Spanish: \emph{tener celos} — to be jealous.

\begin{cousinweb}
\emph{celos} is the plural of \emph{celo} \textquotedblleft{}zeal\textquotedblright{}, which comes from Latin \emph{zēlus}, borrowed from Greek \emph{zēlos} \textquotedblleft{}zeal, eager rivalry, jealousy\textquotedblright{} (EtymDB).
\end{cousinweb}

\begin{grammarlens}[title={What you've built}]
One more word for telling, explaining and describing at B1.
\end{grammarlens}

\subsection*{Your turn}
\begin{itemize}
\raggedright
  \item \emph{Say it:} \emph{los celos}
  \item \emph{Say it:} \emph{los celos}, once more
  \item \emph{Say it:} a sentence of your own with \emph{los celos}
  \item \emph{Recall:} say the Spanish for anger, then the Spanish for to annoy, then say \emph{los celos} again
\end{itemize}

\begin{tcolorbox}[breakable,colback=teal!4,colframe=teal!35!black,title={Before you move on}]
What does los celos mean? (\textquotedblleft{}Jealousy\textquotedblright{}.) Say it once more.
\end{tcolorbox}
\section[celoso]{celoso — jealous}

\label{lesson:ES-C708-celoso}



\begin{quote}
Before the new one: say the Spanish for to annoy, then the Spanish for jealousy.
\end{quote}

\subsection*{You'll want to know: celoso}
\textbf{celoso} — \textquotedblleft{}jealous\textquotedblright{}.

\begin{quote}
Mi novio es muy celoso.
\end{quote}

\begin{quote}
\emph{My boyfriend is very jealous.}
\end{quote}

\emph{Ser celoso} is your character; \emph{estar celoso} is how you feel now.

\begin{cousinweb}
\emph{celoso} is inherited from Late Latin \emph{zēlōsus}, built on \emph{zēlus} \textquotedblleft{}zeal\textquotedblright{} (EtymDB). English \emph{jealous}, through French, and \emph{zealous} both go back to the same Latin family.
\end{cousinweb}

\begin{grammarlens}[title={What you've built}]
One more word for telling, explaining and describing at B1.
\end{grammarlens}

\subsection*{Your turn}
\begin{itemize}
\raggedright
  \item \emph{Say it:} \emph{celoso}
  \item \emph{Say it:} \emph{celoso}, once more
  \item \emph{Say it:} a sentence of your own with \emph{celoso}
  \item \emph{Recall:} say the Spanish for to annoy, then the Spanish for jealousy, then say \emph{celoso} again
\end{itemize}

\begin{tcolorbox}[breakable,colback=teal!4,colframe=teal!35!black,title={Before you move on}]
What does celoso mean? (\textquotedblleft{}Jealous\textquotedblright{}.) Say it once more.
\end{tcolorbox}
\section[la envidia]{la envidia — envy}

\label{lesson:ES-C708-envidia}



\begin{quote}
Before the new one: say the Spanish for jealousy, then the Spanish for jealous.
\end{quote}

\subsection*{You'll want to know: la envidia}
\textbf{la envidia} — \textquotedblleft{}envy\textquotedblright{}.

\begin{quote}
Me da envidia tu viaje.
\end{quote}

\begin{quote}
\emph{I envy you your trip.}
\end{quote}

\emph{Envidia sana} — healthy envy, the friendly kind.

\begin{cousinweb}
\emph{envidia} was borrowed from Latin \emph{invidia} \textquotedblleft{}envy\textquotedblright{}, built on \emph{invidēre} \textquotedblleft{}to eye with ill will\textquotedblright{}, from \emph{in} and \emph{vidēre} \textquotedblleft{}to see\textquotedblright{} (EtymDB).
\end{cousinweb}

\begin{grammarlens}[title={What you've built}]
One more word for telling, explaining and describing at B1.
\end{grammarlens}

\subsection*{Your turn}
\begin{itemize}
\raggedright
  \item \emph{Say it:} \emph{la envidia}
  \item \emph{Say it:} \emph{la envidia}, once more
  \item \emph{Say it:} a sentence of your own with \emph{la envidia}
  \item \emph{Recall:} say the Spanish for jealousy, then the Spanish for jealous, then say \emph{la envidia} again
\end{itemize}

\begin{tcolorbox}[breakable,colback=teal!4,colframe=teal!35!black,title={Before you move on}]
What does la envidia mean? (\textquotedblleft{}Envy\textquotedblright{}.) Say it once more.
\end{tcolorbox}
\section[envidiar]{envidiar — to envy}

\label{lesson:ES-C708-envidiar}



\begin{quote}
Before the new one: say the Spanish for jealous, then the Spanish for envy.
\end{quote}

\subsection*{You'll want to know: envidiar}
\textbf{envidiar} — \textquotedblleft{}to envy\textquotedblright{}.

\begin{quote}
No envidio a los famosos.
\end{quote}

\begin{quote}
\emph{I do not envy famous people.}
\end{quote}

\emph{Envidiar algo a alguien} — to envy someone something.

\begin{cousinweb}
\emph{envidiar} is built on \emph{envidia}, which was borrowed from Latin \emph{invidia} \textquotedblleft{}envy\textquotedblright{}, from \emph{vidēre} \textquotedblleft{}to see\textquotedblright{} (EtymDB).
\end{cousinweb}

\begin{grammarlens}[title={What you've built}]
One more word for telling, explaining and describing at B1.
\end{grammarlens}

\subsection*{Your turn}
\begin{itemize}
\raggedright
  \item \emph{Say it:} \emph{envidiar}
  \item \emph{Say it:} \emph{envidiar}, once more
  \item \emph{Say it:} a sentence of your own with \emph{envidiar}
  \item \emph{Recall:} say the Spanish for jealous, then the Spanish for envy, then say \emph{envidiar} again
\end{itemize}

\begin{tcolorbox}[breakable,colback=teal!4,colframe=teal!35!black,title={Before you move on}]
What does envidiar mean? (\textquotedblleft{}To envy\textquotedblright{}.) Say it once more.
\end{tcolorbox}
\section[el rencor]{el rencor — resentment, grudge}

\label{lesson:ES-C708-rencor}



\begin{quote}
Before the new one: say the Spanish for envy, then the Spanish for to envy.
\end{quote}

\subsection*{You'll want to know: el rencor}
\textbf{el rencor} — \textquotedblleft{}resentment, grudge\textquotedblright{}.

\begin{quote}
No le guardo rencor.
\end{quote}

\begin{quote}
\emph{I do not hold a grudge against him.}
\end{quote}

\emph{Guardar rencor} — to bear a grudge.

\begin{cousinweb}
\emph{rencor} comes from Old Spanish \emph{rancor}, from Latin \emph{rancor} \textquotedblleft{}a rancid smell, a grudge\textquotedblright{} (EtymDB); English \emph{rancour} is the same word.
\end{cousinweb}

\begin{grammarlens}[title={What you've built}]
One more word for telling, explaining and describing at B1.
\end{grammarlens}

\subsection*{Your turn}
\begin{itemize}
\raggedright
  \item \emph{Say it:} \emph{el rencor}
  \item \emph{Say it:} \emph{el rencor}, once more
  \item \emph{Say it:} a sentence of your own with \emph{el rencor}
  \item \emph{Recall:} say the Spanish for envy, then the Spanish for to envy, then say \emph{el rencor} again
\end{itemize}

\begin{tcolorbox}[breakable,colback=teal!4,colframe=teal!35!black,title={Before you move on}]
What does el rencor mean? (\textquotedblleft{}Resentment, grudge\textquotedblright{}.) Say it once more.
\end{tcolorbox}
